% 致谢
\begin{acknowledgement}

    行文至此，这段漫长而细碎的时光，也在此处悄然落笔。

    这四年里，我收获颇丰。来到了我心心念念的上海，见识到了更广阔的天地，也拥有了十八岁时未曾想象过的眼界；
    告别了曾令我痛苦万分的历史与语文，转而投入自己真正热爱的专业领域，在学习与探索中不断成长；
    结识了一群志趣相投的朋友，也始终未曾放下对羽毛球的热爱。
    我愈发明白，成长从来不是独自跋涉的结果，那些平凡却珍贵的日子里，始终有人给予我帮助与陪伴。正因如此，我才能一步步走到今天。借此机会，我想向一路上支持与鼓励过我的人，表达最诚挚的感谢。

    首先，我要感谢我的毕业设计指导老师滕中梅老师。
    滕老师是我本科生涯中对我影响最深的老师之一。
    她扎实深厚的专业素养、出色的教学能力以及认真负责的工作态度，
    使原本应当略显枯燥的数据结构、软件工程、编译原理等课程变得生动而富有魅力，
    也让我能够始终保持浓厚的学习兴趣，并以轻松而充实的状态高质量地完成相关课程的学习。
    十分庆幸也十分感激能够在本科阶段遇到滕老师。

    感谢我的三位室友：袁澍楠、韩淏名和章时豪。大学四年的朝夕相处中，
    你们给予了我许多学习与生活上的帮助与陪伴。
    无论是课程学习中的交流讨论，还是日常生活中的相互照顾；
    无论是面对压力与迷茫时的倾听与鼓励，还是平凡日子里的欢笑与陪伴，
    都让我的大学时光变得更加温暖而充实。很幸运能够与你们相遇，共同度过人生中这段珍贵而难忘的青春岁月。

    最后，我想把最真挚、最深的感谢留给我的爸爸妈妈。感谢你们从小到大给予我无条件的关爱与支持，为我提供了优渥的成长环境和良好的教育条件，让我能够没有后顾之忧地学习、生活，去接触更广阔的世界、追求自己真正热爱的事情。一路走来，我所拥有的一切收获与成绩，都离不开你们长期以来默默的付出与陪伴。你们不仅给予了我物质上的保障，更给予了我面对生活的底气与勇气。能够成为你们的孩子，是我一直以来最幸运的事情。




\rightline{李耀东}
\rightline{宁波}
\rightline{2026 年 5 月 13 日夜}

\end{acknowledgement}
